\subsection{An exact rational positivity check}

A shorter certificate of positivity uses only the first ten
Dirichlet coefficients and rational arithmetic. The sign can be
obtained without an interval evaluation of $\theta(1)$:
$C$ is isomorphic to $y^2=x^5+256$ via
$x=4u$, $y=16(2v+1)$. In
\citet[Proposition 2.2(4)--(6)]{StollYang2003}, the parameter $256=2^8$
has conductor exponent zero at $2$ and
$(256^4-1)/5\equiv4\pmod5$. Their formula gives conductor
$(1-\zeta_5)^2$ and root number $+1$.
Their Lemma 2.1 also verifies the trivial-central-character
normalization of the induced Hilbert form.

\begin{proposition}\label{prop:rationalpositive}
There is an exact rational certificate that $L(A/\Q,1)>1/20$.
\end{proposition}
\begin{proof}
For $p=2,3,7$ and $q=p,p^2$, the map $x\mapsto x^5$ is bijective
on $\F_q$. Thus $\#C(\F_q)=q+1$ and $P_p(T)=1+p^2T^4$.
The bad Euler factor at $5$ is one. Multiplicativity now gives
$a_1=1$ and $a_2=\cdots=a_{10}=0$.
Writing $z_n=c\sqrt n$, equation~\eqref{eq:central} becomes
\[
L(A/\Q,1)=2\sum_{n\geq1}\frac{a_n}{n}z_nK_1(z_n).
\]
The integral representation of $K_1$ and
\[
1+x/2-x^2/8\leq\sqrt{1+x}\leq1+x/2\qquad(x\geq0)
\]
give the global bounds
\[
\begin{split}
K_1(z)&\geq\sqrt{\frac\pi{2z}}e^{-z}
             \left(1+\frac3{8z}-\frac{15}{128z^2}\right),\\
K_1(z)&\leq\sqrt{\frac\pi{2z}}e^{-z}
             \left(1+\frac3{8z}\right).
\end{split}
\]
Indeed the remaining integral is the expectation of
$\sqrt{1+t/(2z)}$ against the density
$e^{-t}t^{1/2}/\Gamma(3/2)$, whose first two moments are $3/2$ and
$15/4$. These inequalities are valid for all positive $z$.

Use $|a_n|\leq d_4(n)\sqrt n$. The rational-arithmetic program
bounds the term at one below and the terms from $11$ through $400$
above using the displayed inequalities. It encloses $\pi$ by
Machin's formula and alternating arctangent sums, each square root
by rational numbers checked by squaring, and each exponential by
positive Taylor sums and a geometric remainder. Thus every
rounding direction is an exact rational inequality.
For the infinite tail note that $d_4(n)\leq n^2$; this follows
multiplicatively from $\binom{e+3}{3}\leq p^{2e}$.
The inequalities $c>5/3$ and $\sqrt{\pi c/2}<2$ give an upper
bound $8n^2e^{-(5/3)\sqrt n}$ for each absolute summand with $n>400$.
This decreases there, so its sum is at most
\[
16e^{-100/3}\sum_{j=0}^{5}
 \frac{5!}{(5-j)!}\,20^{5-j}\left(\frac35\right)^{j+1}.
\]
This last exponential is enclosed by the same positive Taylor sums.
The resulting exact inequalities, retained as executable assertions,
are
\[
2z_1K_1(z_1)>\frac7{10},\qquad
2\sum_{n\geq11}\frac{d_4(n)}{\sqrt n}z_nK_1(z_n)<\frac{13}{20}.
\]
They imply $L(A/\Q,1)>7/10-13/20=1/20$.
\end{proof}
